\chapter{Conclusion}
\section{Lessons Learnt}
	\begin{itemize}
		\item Studied helicopter gearbox topologies, including those used in the A.R.T. program, and 
		compared stage-layout options, gear-pairing strategies, and certification constraints to see how 
		design decisions shape overall geometry and mass.
		\item Observed how varying single operational parameters causes complex changes across the entire gearbox design 
		and that mass does not scale linearly with torque or input power.
		\item Learnt that machine learning models designed for continuous spaces naturally fail to approximate step-functions.
		\item Utilized model-agnostic feature evaluation to rank which operational parameters truly dominate final mass and size.
		\item Learned Romax Concept (Student Edition) in depth: how it works, creating gearbox configurations manually, 
		running simulations, and learning about safety factors and design capabilities.
		\item Gained hands-on experience using Romax Concept (Student Edition) to build, simulate, and validate 
		complex gear arrangements.
		\item Learnt range of machine-learning models (linear regression, random forests, gradient-boosted trees, 
		and neural networks) evaluating their strengths and limits for predicting size, weight, and performance.
		\item Learnt how to distribute points in a N-dimensional space and why completely random sampling leaves blind spots 
		for the AI, and how quasi-random Sobol Sequences solve this by spacing out training points uniformly.
	\end{itemize}

\section{Conclusion}
	The workflow introduced in this work replaces the traditional, time-consuming, and skill-dependent gearbox design cycle 
	with a repeatable, data-driven process that still respects every engineering constraint. The sizing engine, implemented 
	in Python, accepts 31 design, material, and operating variables and returns the full geometric description of a 
	dual-stage helicopter main gearbox (spiral-bevel input and planetary output).  

	The design parameters produced by the Python design script were manually entered into the Romax Concept (Student Edition). 
	The software then returned the gearbox weight, the stress-analysis results, the actual gear ratio, and the safety factor. 
	The Python output predicts a gear mass of 364 kg (excluding the casing), which is within <2\% of the 357.8 kg mass reported by the 
	Romax Concept (Student Edition). 
	A sensitivity study identified gear modules, pressure and helix angles, and planetary configuration as the most influential 
	parameters for weight. These results informed the sampling strategy for the training and highlighted the location of the 
	dominant design trade-offs. 

	A Sobol-sequence generated tens of thousands of distinct configurations to train two surrogate AI models. 
	A first-stage “Sizer” (HistGradientBoosting) predicts the continuous gear module; the output is then snapped 
	to the nearest ISO standard module. A second-stage “Weigher” (also HistGradientBoosting) predicts total mass 
	and volume from the snapped module and the remaining design variables.  

	The cascade approach eliminates the discontinuity that appears if a single model is trained on the 
	raw, continuous outputs. The final surrogate achieves ($R^{2}=0.984$) and a mean absolute error of 15.36 kg. 

	Future work will extend the database to cover a wider range of transmission configurations and 
	validate the predicted designs through prototype testing. The framework presented here provides a 
	scalable foundation for rapid, physics-based design of high-performance gearboxes in aerospace applications. 